% CP-VI-0023
% target_chapter: VI/32 — Tangent Spaces and Local Geometry
% source_unit_id: IMP-DL-721195CCB5-PSR001753-8BBBA0F3
% source: Downloads/theory-of-algebraic-geometry-11.tex L1753-L2613
% solution_provenance: SOURCE_IMPLICIT_SOLUTION

\subsection*{Problem 6: Relative Differentials}

\subsection*{(a) Construction and Universal Property of \(\Omega_{B/A}\)}

Let \(A\) be a ring, let \(B\) be an \(A\)-algebra, and let \(M\) be a
\(B\)-module.

An \(A\)-derivation of \(B\) into \(M\) is a map
\[
d:B\to M
\]
such that:
\[
d(b+b')=d(b)+d(b'),
\]
\[
d(bb')=b\,d(b')+b'\,d(b),
\]
and
\[
d(a)=0
\qquad
\text{for all }a\in A.
\]

Let
\[
\mu:B\otimes_A B\longrightarrow B
\]
be the multiplication map
\[
\mu(b\otimes b')=bb'.
\]

Let
\[
I=\ker(\mu).
\]

Define
\[
\Omega_{B/A}=I/I^2.
\]

For \(b\in B\), define
\[
db=1\otimes b-b\otimes1.
\]

Since
\[
\mu(1\otimes b-b\otimes1)=b-b=0,
\]
we have
\[
db\in I.
\]

Thus \(db\) defines an element of \(I/I^2\).

We check that
\[
d:B\to I/I^2
\]
is an \(A\)-derivation.

First,
\[
d(b+b')
=
1\otimes(b+b')-(b+b')\otimes1.
\]

Therefore
\[
d(b+b')=db+db'.
\]

Second,
\[
d(bb')
=
1\otimes bb'-bb'\otimes1.
\]

Modulo \(I^2\), one has
\[
d(bb')=b\,db'+b'\,db.
\]

Indeed, this identity is exactly the image in \(I/I^2\) of the standard
identity expressing the Leibniz rule.

Finally, if \(a\in A\), then
\[
da=1\otimes a-a\otimes1=0
\]
because the tensor product is taken over \(A\).

So \(d\) is an \(A\)-derivation.

Now let
\[
d':B\to M
\]
be any \(A\)-derivation.

Define a map
\[
\varphi:B\otimes_A B\to M
\]
by
\[
\varphi(b\otimes b')=b\,d'(b').
\]

The derivation rule implies that \(\varphi\) sends \(I^2\) to zero after
restriction to \(I\). Hence \(\varphi|_I\) factors through \(I/I^2\).

Therefore there exists a \(B\)-linear map
\[
f:I/I^2\to M
\]
such that
\[
f(db)=d'(b).
\]

Thus
\[
d'=f\circ d.
\]

The map \(f\) is unique because \(\Omega_{B/A}=I/I^2\) is generated by the
elements \(db\).

Therefore \(\Omega_{B/A}\), together with \(d:B\to\Omega_{B/A}\), satisfies
the universal property:

for every \(A\)-derivation \(d':B\to M\), there exists a unique \(B\)-linear
map
\[
f:\Omega_{B/A}\to M
\]
such that
\[
d'=f\circ d.
\]

Hence
\[
\boxed{
	\Omega_{B/A}=I/I^2.
}
\]

\subsection*{(b) Computation of \(\Omega_{k[x_1,\ldots,x_n]/k}\)}

Let
\[
A=k[x_1,\ldots,x_n].
\]

We claim that
\[
\Omega_{A/k}
\cong
A\,dx_1\oplus\cdots\oplus A\,dx_n.
\]

Define
\[
d:A\to \bigoplus_{i=1}^n A\,dx_i
\]
by
\[
d(f)=\sum_{i=1}^n
\frac{\partial f}{\partial x_i}dx_i.
\]

This map is \(k\)-linear and satisfies
\[
d(fg)=f\,dg+g\,df
\]
because ordinary partial differentiation satisfies the product rule.

Also,
\[
d(c)=0
\qquad
(c\in k).
\]

Therefore \(d\) is a \(k\)-derivation.

Now let
\[
D:A\to M
\]
be any \(k\)-derivation into an \(A\)-module \(M\).

Since \(A\) is a polynomial ring, \(D\) is determined by the values
\[
D(x_1),\ldots,D(x_n).
\]

Define an \(A\)-linear map
\[
\varphi:
A\,dx_1\oplus\cdots\oplus A\,dx_n
\longrightarrow M
\]
by
\[
\varphi(dx_i)=D(x_i).
\]

Then for every polynomial \(f\),
\[
\varphi(df)
=
\sum_{i=1}^n
\frac{\partial f}{\partial x_i}D(x_i).
\]

By the usual formula for derivations on a polynomial ring,
\[
D(f)
=
\sum_{i=1}^n
\frac{\partial f}{\partial x_i}D(x_i).
\]

Thus
\[
D=\varphi\circ d.
\]

The map \(\varphi\) is unique because \(dx_1,\ldots,dx_n\) form a basis.

Hence
\[
\boxed{
	\Omega_{k[x_1,\ldots,x_n]/k}
	\cong
	\bigoplus_{i=1}^n k[x_1,\ldots,x_n]\,dx_i.
}
\]

So \(\Omega_{A/k}\) is the free \(A\)-module generated by
\[
dx_1,\ldots,dx_n.
\]

Moreover, for every polynomial \(f\in A\),
\[
\boxed{
	df=\sum_{i=1}^n
	\frac{\partial f}{\partial x_i}dx_i.
}
\]

\subsection*{(c) The Sheaf of Differentials on Affine Schemes}

Let
\[
X=\operatorname{Spec}B,
\qquad
Y=\operatorname{Spec}A,
\]
where \(B\) is an \(A\)-algebra.

Then
\[
X\times_Y X
=
\operatorname{Spec}(B\otimes_A B).
\]

The diagonal morphism
\[
\Delta:X\to X\times_Y X
\]
corresponds to the multiplication map
\[
B\otimes_A B\to B,
\qquad
b\otimes b'\mapsto bb'.
\]

Let
\[
I=\ker(B\otimes_A B\to B).
\]

The conormal sheaf of the diagonal closed immersion is
\[
I/I^2
\]
sheafified on \(X\).

By part (a),
\[
I/I^2=\Omega_{B/A}.
\]

Therefore
\[
\Omega_{X/Y}
\cong
\widetilde{\Omega_{B/A}}.
\]

Hence
\[
\boxed{
	\Omega_{\operatorname{Spec}B/\operatorname{Spec}A}
	\cong
	\widetilde{\Omega_{B/A}}.
}
\]

\subsection*{(d) Differentials on Affine Space}

Let
\[
X=\mathbb A_k^n
=
\operatorname{Spec}k[x_1,\ldots,x_n].
\]

By part (b),
\[
\Omega_{k[x_1,\ldots,x_n]/k}
\cong
\bigoplus_{i=1}^n k[x_1,\ldots,x_n]\,dx_i.
\]

Sheafifying this module on \(\mathbb A_k^n\), we obtain
\[
\Omega_{\mathbb A_k^n/\operatorname{Spec}k}
\cong
\bigoplus_{i=1}^n
\mathcal O_{\mathbb A_k^n}\,dx_i.
\]

Therefore
\[
\boxed{
	\Omega_{\mathbb A_k^n/\operatorname{Spec}k}
	\cong
	\mathcal O_{\mathbb A_k^n}^{\oplus n}.
}
\]

Thus \(\Omega_{\mathbb A_k^n/\operatorname{Spec}k}\) is a free
\(\mathcal O_{\mathbb A_k^n}\)-module of rank \(n\), generated by
\[
dx_1,\ldots,dx_n.
\]

\section{Theory: Affine Covers, \v{C}ech Cohomology, Plane Curves, and Relative Differentials}

\subsection{Affine Schemes and Vanishing of Higher Cohomology}

A fundamental theorem in algebraic geometry says that if $X$ is affine and
$\mathcal F$ is quasi-coherent, then
\[
H^p(X,\mathcal F)=0
\qquad
\text{for all }p>0.
\]

In particular, for the structure sheaf,
\[
H^1(X,\mathcal O_X)=0
\]
whenever $X$ is affine.

Therefore, to prove that a scheme $U$ is not affine, it is enough to find
\[
H^1(U,\mathcal O_U)\neq 0.
\]

For example, for the punctured affine plane
\[
U=\mathbb A_k^2\setminus\{0\},
\]
we use the cover
\[
U=D(x)\cup D(y).
\]

The overlap is
\[
D(x)\cap D(y)=D(xy).
\]

The \v{C}ech computation gives
\[
H^1(U,\mathcal O_U)
\cong
\frac{k[x,y]_{xy}}{k[x,y]_x+k[x,y]_y}.
\]

Now
\[
k[x,y]_{xy}=k[x,y,x^{-1},y^{-1}].
\]

The monomials which survive in the quotient are exactly those of the form
\[
x^iy^j
\qquad
(i<0,\ j<0).
\]

Hence
\[
H^1(U,\mathcal O_U)
\cong
\bigoplus_{i<0,\ j<0}kx^iy^j.
\]

This is infinite-dimensional over $k$, and in particular nonzero. Therefore
\[
\mathbb A_k^2\setminus\{0\}
\]
is not affine.

The principle is:
\[
\boxed{
	H^1(U,\mathcal O_U)\neq0
	\quad\Longrightarrow\quad
	U\text{ is not affine}.
}
\]

\subsection{\v{C}ech Cohomology on a Two-Open Affine Cover}

Suppose
\[
X=U\cup V
\]
where $U$, $V$, and $U\cap V$ are affine.

Then the \v{C}ech complex for $\mathcal O_X$ is
\[
0
\longrightarrow
\mathcal O_X(X)
\longrightarrow
\mathcal O_X(U)\oplus\mathcal O_X(V)
\longrightarrow
\mathcal O_X(U\cap V)
\longrightarrow 0.
\]

The map is
\[
(s_U,s_V)
\longmapsto
s_V|_{U\cap V}-s_U|_{U\cap V}.
\]

Thus
\[
H^0(X,\mathcal O_X)
=
\ker\left(
\mathcal O_X(U)\oplus\mathcal O_X(V)
\to
\mathcal O_X(U\cap V)
\right),
\]
and
\[
H^1(X,\mathcal O_X)
=
\operatorname{coker}\left(
\mathcal O_X(U)\oplus\mathcal O_X(V)
\to
\mathcal O_X(U\cap V)
\right).
\]

Therefore $H^0$ consists of pairs of regular functions on $U$ and $V$ that
agree on the overlap, while $H^1$ measures regular functions on the overlap
which cannot be written as a difference of functions coming from $U$ and $V$.

This is the reason why quotients such as
\[
\frac{k[x,y]_{xy}}{k[x,y]_x+k[x,y]_y}
\]
naturally appear in \v{C}ech cohomology.

\subsection{Projective Plane Curves and the Genus Formula}

Let
\[
X\subseteq \mathbb P_k^2
\]
be a plane curve defined by one homogeneous polynomial
\[
f(x_0,x_1,x_2)=0
\]
of degree $d$.

Assume
\[
(1,0,0)\notin X.
\]

Then $X$ is covered by
\[
U=X\cap D_+(x_1),
\qquad
V=X\cap D_+(x_2).
\]

Indeed, the complement of
\[
D_+(x_1)\cup D_+(x_2)
\]
in $\mathbb P_k^2$ is the single point
\[
(1,0,0),
\]
which does not belong to $X$.

Thus
\[
X=U\cup V.
\]

The sets $U$, $V$, and $U\cap V$ are affine. Therefore the two-open
\v{C}ech complex computes $H^0(X,\mathcal O_X)$ and
$H^1(X,\mathcal O_X)$.

For a projective plane curve, global regular functions are constants, so
\[
H^0(X,\mathcal O_X)\cong k.
\]

Hence
\[
\dim_k H^0(X,\mathcal O_X)=1.
\]

The first cohomology group $H^1(X,\mathcal O_X)$ measures the failure of
regular functions on the overlap $U\cap V$ to be represented by functions on
$U$ and $V$.

For a plane curve of degree $d$, the explicit \v{C}ech computation gives
\[
\dim_k H^1(X,\mathcal O_X)
=
\frac{(d-1)(d-2)}2.
\]

If $X$ is nonsingular, then this dimension equals the genus:
\[
g(X)=\dim_k H^1(X,\mathcal O_X).
\]

Therefore a nonsingular plane curve of degree $d$ has genus
\[
\boxed{
	g=\frac{(d-1)(d-2)}2.
}
\]

For example:
\[
d=1\quad\Rightarrow\quad g=0,
\]
\[
d=2\quad\Rightarrow\quad g=0,
\]
\[
d=3\quad\Rightarrow\quad g=1,
\]
\[
d=4\quad\Rightarrow\quad g=3.
\]

Thus lines and conics have genus $0$, smooth cubics have genus $1$, and
smooth quartics have genus $3$.

\subsection{K\"ahler Differentials}

Let
\[
A\to B
\]
be a ring homomorphism.

The module of relative differentials
\[
\Omega_{B/A}
\]
is the universal object representing $A$-derivations from $B$.

An $A$-derivation into a $B$-module $M$ is a map
\[
D:B\to M
\]
such that
\[
D(b+b')=D(b)+D(b'),
\]
\[
D(bb')=bD(b')+b'D(b),
\]
and
\[
D(a)=0
\qquad
(a\in A).
\]

The universal property of $\Omega_{B/A}$ says that there exists an
$A$-derivation
\[
d:B\to\Omega_{B/A}
\]
such that every $A$-derivation
\[
D:B\to M
\]
factors uniquely through $d$.

That is, there is a unique $B$-linear map
\[
\varphi:\Omega_{B/A}\to M
\]
such that
\[
D=\varphi\circ d.
\]

Equivalently,
\[
\operatorname{Der}_A(B,M)
\cong
\operatorname{Hom}_B(\Omega_{B/A},M).
\]

Thus $\Omega_{B/A}$ represents the functor of $A$-derivations.

\subsection{Construction of Differentials Using the Diagonal}

The module $\Omega_{B/A}$ can be constructed as
\[
\Omega_{B/A}=I/I^2,
\]
where
\[
I=\ker(B\otimes_A B\to B),
\]
and the map
\[
B\otimes_A B\to B
\]
is multiplication:
\[
b\otimes b'\longmapsto bb'.
\]

The universal differential is
\[
db=1\otimes b-b\otimes1.
\]

Indeed,
\[
1\otimes b-b\otimes1\in I
\]
because multiplication sends it to
\[
b-b=0.
\]

The quotient $I/I^2$ keeps only first-order information. Products of elements
of $I$ are second-order infinitesimals, and these are killed in $I/I^2$.

Geometrically, this construction comes from the diagonal morphism
\[
\Delta:X\to X\times_Y X.
\]

When
\[
X=\operatorname{Spec}B,
\qquad
Y=\operatorname{Spec}A,
\]
we have
\[
X\times_Y X=\operatorname{Spec}(B\otimes_A B).
\]

The diagonal corresponds to the multiplication map
\[
B\otimes_A B\to B.
\]

The ideal $I$ defines the diagonal, and
\[
I/I^2
\]
is the conormal module of the diagonal.

Thus differentials measure first-order infinitesimal displacement away from
the diagonal.

\subsection{Differentials on Affine Space}

Let
\[
A=k[x_1,\ldots,x_n].
\]

Then
\[
\Omega_{A/k}
\cong
A\,dx_1\oplus\cdots\oplus A\,dx_n.
\]

Thus $\Omega_{A/k}$ is a free $A$-module of rank $n$.

For every polynomial $f\in A$,
\[
df
=
\sum_{i=1}^n
\frac{\partial f}{\partial x_i}dx_i.
\]

This is the algebraic analogue of the total differential from calculus.

After sheafifying on
\[
\mathbb A_k^n=\operatorname{Spec}k[x_1,\ldots,x_n],
\]
we get
\[
\Omega_{\mathbb A_k^n/\operatorname{Spec}k}
\cong
\mathcal O_{\mathbb A_k^n}^{\oplus n}.
\]

Therefore
\[
\Omega_{\mathbb A_k^n/\operatorname{Spec}k}
\]
is a free $\mathcal O_{\mathbb A_k^n}$-module of rank $n$, generated by
\[
dx_1,\ldots,dx_n.
\]

\subsection{Conceptual Summary}

The problems in this group are linked by the same principle:

\[
\boxed{
	\text{local algebra controls global geometry.}
}
\]

For the punctured affine plane, the localizations
\[
k[x,y]_x,
\qquad
k[x,y]_y,
\qquad
k[x,y]_{xy}
\]
detect non-affineness through the nonzero cohomology group
\[
H^1(U,\mathcal O_U).
\]

For plane curves, affine charts and their overlap compute
\[
H^0(X,\mathcal O_X)
\qquad\text{and}\qquad
H^1(X,\mathcal O_X),
\]
leading to the genus formula
\[
g=\frac{(d-1)(d-2)}2.
\]

For relative differentials, the algebra
\[
B\otimes_A B
\]
and the diagonal ideal
\[
I=\ker(B\otimes_A B\to B)
\]
encode infinitesimal geometry through
\[
\Omega_{B/A}=I/I^2.
\]

Thus these problems introduce three central tools:
\[
\boxed{
	\text{\v{C}ech cohomology, affine covers, and K\"ahler differentials.}
}
\]
\section{Examples: Affine Covers, \v{C}ech Cohomology, Plane Curves, and Differentials}

\subsection{Example 1: A Basic \v{C}ech Computation on Two Opens}

Let
\[
X=U\cup V.
\]

For the structure sheaf, the \v{C}ech complex is
\[
0
\longrightarrow
\mathcal O_X(X)
\longrightarrow
\mathcal O_X(U)\oplus\mathcal O_X(V)
\longrightarrow
\mathcal O_X(U\cap V)
\longrightarrow 0.
\]

The middle map is
\[
(f,g)
\longmapsto
g|_{U\cap V}-f|_{U\cap V}.
\]

Therefore
\[
H^0(X,\mathcal O_X)
=
\{(f,g)\mid f|_{U\cap V}=g|_{U\cap V}\},
\]
and
\[
H^1(X,\mathcal O_X)
=
\frac{\mathcal O_X(U\cap V)}
{\mathcal O_X(U)+\mathcal O_X(V)}.
\]

Thus $H^1$ measures functions on the overlap which cannot be produced from
functions on the two pieces.
